% Generated draft: requires mathematical and source review.
\documentclass[11pt]{article}
\usepackage[margin=25mm]{geometry}
\usepackage{fontspec}
\setmainfont{Latin Modern Roman}
\usepackage{amsmath,amssymb,mathtools,unicode-math}
\setmathfont{Latin Modern Math}
\usepackage{enumitem,needspace,xurl,hyperref}
\setlength{\parindent}{0pt}
\setlength{\parskip}{5pt}
\setlength{\emergencystretch}{3em}
\newcommand{\N}{\mathbb{N}}
\newcommand{\Z}{\mathbb{Z}}
\newcommand{\Q}{\mathbb{Q}}
\newcommand{\R}{\mathbb{R}}
\newcommand{\C}{\mathbb{C}}
\newcommand{\F}{\mathbb{F}}
\newcommand{\A}{\mathbb{A}}
\newcommand{\PP}{\mathbb{P}}
\newcommand{\E}{\mathbb{E}}
\newcommand{\eps}{\varepsilon}
\newcommand{\dd}{\,\mathrm{d}}
\begin{document}
\section*{UnsolvedMath ID: 30004018}
\textbf{Source problem code:} catalogue entry
\subsection{Definitions and mathematical statement}
The supplied record concerns a module denoted by \(K_{m\mid n}\) and Duflo--Serganova functors denoted by \(\operatorname{DS}_x\). Its background identifies the intended setting as category \(\mathcal O\), but does not specify the underlying algebra, the precise category, the construction of \(K_{m\mid n}\), its acting algebra or coefficient field, or the sets \(X_k\).

For each index \(k\) satisfying the supplied restriction \(k\leq\min(m,n)\), define \(K^k_{m\mid n}\), in the source's intended construction, to be the subobject generated by all objects \(M\) satisfying
\[
\operatorname{DS}_x(M)\cong 0\qquad\text{for every }x\in X_k.
\]
Here annihilation means that the functor has zero value on the object itself. The record does not specify how objects give generators of \(K_{m\mid n}\), or whether generation means linear span, subgroup generation, or generation as a module; these are unresolved parts of this definition.

For a module \(V\), its socle \(\operatorname{soc}(V)\) is the sum of its simple submodules. Under the convention that the ascending socle filtration starts at zero, its terms are defined by
\[
\operatorname{soc}_0(V)=0,\qquad
\operatorname{soc}_{j+1}(V)=\pi_j^{-1}\!\left(\operatorname{soc}\!\left(V/\operatorname{soc}_j(V)\right)\right)
\quad(j\geq 0),
\]
where \(\pi_j:V\to V/\operatorname{soc}_j(V)\) is the quotient map. This explains a standard convention; the supplied record does not confirm its indexing convention or specify the acting algebra whose simple submodules are meant.

The proposed equality is
\[
\operatorname{soc}_k K_{m\mid n}=K^k_{m\mid n}
\qquad\text{for every intended admissible }k\leq\min(m,n).
\]
The domains of \(m,n,k\), the lower bound on \(k\), and inclusion or exclusion of \(k=0\) are not supplied. Thus the displayed formula preserves the source's conjecture, but a fully specified mathematical assertion requires the missing constructions, parameter domains, and filtration convention. In particular, no meaning of \(X_k\) as a rank stratum is assumed here.
\subsection{Short English statement}
The conjecture asks whether each indicated socle layer of \(K_{m\mid n}\) equals the subobject generated by objects killed by every Duflo--Serganova functor indexed by \(X_k\). The supplied background places the full question in category \(\mathcal O\). Missing definitions and indexing conventions prevent a fully self-contained formulation. The supplied dated assessment labels the problem partially solved and distinguishes a result for a finite-dimensional direct summand from the unresolved full conjecture; that assessment is not independently verified here.
\subsection{Sources}
\begin{itemize}
\item Local UnsolvedMath registry: supplied complete record with permanent ID \(30004018\), title \emph{Socle Filtrations from Duflo--Serganova Kernels}, exact source problem code \texttt{catalogue entry}, in the Oberwolfach Reports Open Problems collection. The record gives proposed year 2018 and names no individual proposer.
\item Supplied report citation: \emph{Enveloping Algebras and Geometric Representation Theory} (2019), Oberwolfach Reports, pp.\ 3081--3146, DOI \texttt{10.4171/owr/2018/52}. \url{https://doi.org/10.4171/owr/2018/52}. The supplied background describes Oberwolfach Reports as documenting mathematical workshops and their open-problem discussions.
\item Supplied literature source: \emph{Integrable sl-infinity Modules and Category O}. \url{https://arxiv.org/abs/1712.00664}.
\item Supplied literature source: \emph{The Duflo--Serganova Functor, Vingt Ans Apres}. \url{https://par.nsf.gov/servlets/purl/10512366}.
\item Supplied literature source: \emph{Oberwolfach Report 52/2018}. \url{https://ems.press/content/serial-article-files/46773}.
\item Literature assessment date: \texttt{2026-08-21}. This is supplied dated curation derived from the supplied OWR curation, not an independently verified literature assessment. It labels the problem partially solved and states that the socle filtration and Duflo--Serganova maps are rigorously described, that the tensor-filtration conjecture is proved for the finite-dimensional direct summand, and that the full Conjecture 9 for category \(\mathcal O\) remains unresolved. These source-specific claims and their current status require independent verification.
\item The registry labels the statement \texttt{corrected\_verified}, with a supplied explanation that mathematical content was preserved while presentation or omitted definitions were repaired. This is a registry qualification, not mathematical verification performed for this draft.
\end{itemize}
\end{document}
